\begin{figure}[H]
Jeżeli typ pola nie jest podany to jest to string.
\begin{verbatim}
{
  "id": "UUID klienta",
  "name": "Imię i nazwisko lub nazwa firmy",
  "email": "Adres e-mail klienta"
}
\end{verbatim}
\caption{Obiekt Client}
\label{schema:client}
\end{figure}

\begin{figure}[H]
Jeżeli typ pola nie jest podany to jest to string.
\begin{verbatim}
{
  "id": "UUID rezerwacji",
  "clientId": "UUID klienta",
  "items": [
    "Zrzut stanu i cen pozycji koszyka -- lista obiektów OrderItem"
  ],
  "createTimestamp": "Znacznik czasu utworzenia",
  "lastModifiedTimestamp": "Znacznik ostatniej modyfikacji"
}
\end{verbatim}
\caption{Obiekt Reservation}
\label{schema:reservation}
\end{figure}

\begin{figure}[H]
Jeżeli typ pola nie jest podany to jest to string.
\begin{verbatim}
{
  "lineId": "UUID pozycji koszyka",
  "offerSnapshot": "zrzut oferty w momencie rezerwacji -- obiekt Offer",
  "quantity": "Liczba zarezerwowanych sztuk -- liczba całkowita",
  "timestamp": "Czas utworzenia pozycji (UTC, ISO-8601)",
  "status": "RESERVED / PAID / CANCELLED",
  "participant": {
    "name": "Imię i nazwisko uczestnika",
    "email": "Opcjonalny adres e-mail"
  }
}
\end{verbatim}
\caption{Obiekt OrderItem}
\label{schema:order-item}
\end{figure}

Pole \texttt{participant} jest opcjonalne i może mieć wartość \texttt{null}.
Nowa pozycja ma status \texttt{RESERVED}. Statusy \texttt{PAID} i \texttt{CANCELLED}
są przewidziane dla przyszłych modułów płatności i anulowania; obecne API nie
pozwala ich ustawiać. Oferta w \texttt{offerSnapshot} jest niezmiennym zrzutem
z chwili rezerwacji, przed odjęciem zapasu.
